\begin{claim}{C.80}{Proved}{Every based sphere family fills in the one-dimensional proper disk centralizer}
Let $O=\operatorname{int}\overline D$ be the actual
open disk, and let $Q\leq\operatorname{Homeo}(O)$
be closed, locally compact, free and proper, with
$\dim Q=1$. Retain the actual group and topology
\[
 \mathscr C_Q=
 \{h\in\operatorname{Homeo}(\overline D):
    h|_{\partial\overline D}=\id,
    \ hq=qh\text{ on }O\text{ for every }q\in Q\},
                                                               \tag{C.80.1}
\]
using compact-open convergence of maps and inverses.
For every $n\geq1$ and every continuous based map
$f:(S^n,*)\to(\mathscr C_Q,1)$, there is a continuous
map $H:[0,1]\times S^n\to\mathscr C_Q$ satisfying
\[
 \begin{gathered}
 H_{0,\xi}=1,\qquad H_{1,\xi}=f_\xi,\qquad H_{r,*}=1.\\
 \text{Consequently }\quad
 \pi_n(\mathscr C_Q,1)=0\quad(n\geq1).
 \end{gathered}
                                                               \tag{C.80.2}
\]
Both evaluation families $H_{r,\xi}(x)$ and
$H_{r,\xi}^{-1}(x)$ are jointly continuous on the
original closed disk. Every original operator is
retained. This proves the stated based homotopy
groups, without asserting contractibility of the
whole group or describing its components.
\end{claim}

\begin{proof}
\emph{The actual source group and the sphere parameter.}
Put $\Sigma=S^n$ and
$P=f(\Sigma)\cup f(\Sigma)^{-1}$. This is a compact
path-connected family containing $1$. Use
$G=\langle P\rangle$ and its closure $L$ in the
actual group (C.80.1). The latter is a closed
subgroup of the Polish closed-disk homeomorphism
group: the boundary equations and all commutator
equations on $O$ are closed under convergence of
maps and inverses. Thus $L$ has the complete
compatible group topology used in C.76--C.79.

C.33's actual principal coordinate, with its
inverse, has $Q^0$ acting by real translations
in the first coordinate. Hence every $k\in L$
has the form
\[
 \begin{split}
 k(s,t)&=(s+u_k(t),\beta_k(t)),\\
 k^{-1}(s,t)&=
    (s-u_k(\beta_k^{-1}(t)),\beta_k^{-1}(t)).
 \end{split}                                                   \tag{C.80.3}
\]
The base maps are increasing. This holds first
for $f_\xi$, since a path in $\Sigma$ from $*$
to $\xi$ gives an actual based path of maps,
and then for $G$ and its closure. Every finite
word in $P$ has a based path within the same
finite power of $P$: choose a parameter path
for each letter and multiply the resulting
paths, with the indicated inverses. A jointly
continuous choice of those paths is unnecessary.

Let $A=\{t:\beta_{f_\xi}(t)=t\text{ for all }\xi\}$.
It is closed. Every complementary component
$I$ is a transitive base interval for $G$.
Indeed a based path moving a point and its
inverse give a neighborhood in its base orbit;
the orbit is connected because the increasing
union $\bigcup_m P^m(t)$ is connected. A finite
endpoint of that interval orbit is fixed by
all its increasing base maps. This proves
the stated component equality, exactly as
in C.75, for this actual sphere-generated $G$.
Every $k\in L$ fixes $A$ pointwise on the base
and preserves each $I$.

The endpoint characters are the actual evaluations
$c_a(k)=u_k(a)$, $a\in A$. They are continuous
characters by the composition formula (C.75.6).
All other source facts needed from C.75--C.77
have the same precise hypotheses here. For an
abelian effective base image, its closure is
abelian, transitive and free; the verified
evaluation homeomorphism identifies the actual
image with real translations, as in C.77.
This gives an increasing $b_I:I\to\mathbb R$
and the continuous character
\[
 d_I(k)=b_I(\beta_k(t_0))-b_I(t_0),\qquad
 b_I\beta_k(t)=b_I(t)+d_I(k).
                                                               \tag{C.80.4}
\]
For a nonabelian base image, $[G,G]$ is path
connected: each commutator and each finite
product of commutators has a based path from
the finite-word paths above. Its nontrivial
base image is normal. Its common fixed set
is $G$-invariant and therefore empty by base
transitivity. Its interval orbits are open
by the same path argument. They partition
the connected interval $I$, so it is transitive
there. All endpoint characters vanish on
$[G,G]$. Thus every scalar character/base
comparison used below is exactly the transitive
action in (C.77.3), with its stated kernel.

\emph{Short increments with the complete original topology.}
Give $\Sigma$ its round geodesic metric $d$.
For $0<\delta<\pi$ put
\[
 J_\delta=
 \{f_\xi f_\eta^{-1}:d(\xi,\eta)\leq\delta\},
 \qquad E_\delta=J_\delta^{15}.
                                                               \tag{C.80.5}
\]
The pair space in this formula is compact and
path connected. The unique shortest geodesic
from $\xi$ to $\eta$ contracts its first entry
to the second, continuously on this pair space
and without increasing their distance. The
diagonal is a copy of the path-connected sphere.
Consequently $J_\delta$ is compact, symmetric,
path connected, and contains $1$. Uniform
continuity of $f$ and of multiplication and
inversion on its compact image proves
$J_\delta\to\{1\}$ in the actual topology as
$\delta\downarrow0$.

For any endpoint-character vector, or the
vector consisting of $d_I$ and all finite
endpoint characters, there are at most three
coordinates. Its actual image $V$ is a real
vector subspace. To check this before closure,
apply the five-letter additive path identity
(C.50.3) to a finite-word based path ending
at any specified $g\in G$. This puts every
fraction of its vector in the actual image.
Addition, negatives and integer multiples
then give the whole real span. Applying the
same identity to a path from $*$ to $\xi$
within $\Sigma$ also gives
\[
 r a(f_\xi)\in a(P^5),\qquad
 a(E_\delta)\text{ contains a neighborhood of }0
                          \text{ in }V.
                                                               \tag{C.80.6}
\]
For the second assertion, subdivide a parameter
path into pieces of length less than $\delta$.
Telescoping shows that $a(J_\delta)$ spans $V$.
Choose at most three basis vectors in this
compact symmetric path-connected image. The
same five-letter identity supplies the full
segment from minus to plus each basis vector;
their sum contains a neighborhood of zero.
This uses at most fifteen factors. Every such
character extends continuously to $L$, and
closedness of $V$ gives $a(L)=V$.

Choose the successive $\delta_m$ in (C.80.5)
as in C.77--C.78, so multiplication by the
next compact increment changes every preceding
compact product by less than $2^{-m}$ in a
complete compatible metric. Starting with
$K_0=P^5$, the resulting nested families
$K_m$ lie in one compact path-connected limit
family $K\subset L$. The product parameter
space is a product of compact path-connected
spaces, and uniform Cauchy convergence gives
its continuous map into $L$. These statements
include convergence of the inverse maps.
One such $K$ serves every strip.

\emph{The exact relative parameter conditions.}
Use the compact metric parameter space
\[
 T=[0,1]\times\Sigma,\qquad
 D=(\{0,1\}\times\Sigma)\cup([0,1]\times\{*\}).
                                                               \tag{C.80.7}
\]
On $D$ the prescribed source labels are $1$
at $r=0$, $f_\xi$ at $r=1$, and $1$ at $\xi=*$.
They agree on all intersections because $f_*=1$.
They are continuous on $D$ and belong to $K_0$.

For every nonabelian strip with nonzero endpoint
vector $c_I$, form
\[
 F_m(p)=\{k\in K_m:c_I(k)=r c_I(f_\xi)\},
 \qquad p=(r,\xi)\in T.
                                                               \tag{C.80.8}
\]
These compact fibers are nonempty by (C.80.6).
Use exactly C.78's embedding $\chi:K\to\ell^2$
and compact convex sets $\mathcal P(F)$, with
their continuous positive averages $\Lambda_q$.
The actual character correction from $p$ to
$p'$ is $r'c_I(f_{\xi'})-r c_I(f_\xi)$. By
(C.80.6), an arbitrarily small next increment
realizes it and corrects every label in the
old fiber at once. The finite-convex-combination
argument (C.78.8) gives the interlaced lower-limit
inclusion for these Hilbert convex sets.

The relative selection proof of C.63 and C.78
applies on $T$, with $D$ as its closed relative
subset. Explicitly, replace the allowed sets
on $D$ by the continuous singleton labels
$\chi(1)$ or $\chi(f_\xi)$ just specified.
At a point outside $D$ a neighborhood avoids
$D$; at a point of $D$ the old lower-limit
inclusion and continuity of these singleton
values give the relative inclusion. The
locally finite partitions of unity and
summable successive approximations used in
C.63 require a metrizable paracompact domain,
not a rectangular parameter boundary. Both
$T$ and $T\times I$ have that property.
Their limits lie in the closed limiting
compact-family convex sets, and their errors
to the prescribed singletons tend to zero.
We therefore obtain continuous positive
averages $q_p\in\mathcal P(F_\infty(p))$ with
\[
 q_{0,\xi}=\chi(1),\qquad q_{1,\xi}=\chi(f_\xi),
                       \qquad q_{r,*}=\chi(1).
                                                               \tag{C.80.9}
\]

\emph{All base prescriptions and their compact realizations.}
For an abelian base strip use
$B_p(t)=b_I^{-1}(b_I(t)+r d_I(f_\xi))$.
For a nonabelian base strip with nonzero
endpoint vector use
$B_p(t)=\Lambda_{q_p}(k\mapsto\beta_k(t))$.
In the rank-two case the average uses the
full vector in (C.80.8), so both scalar
endpoint selections below have the same base.
For a nonabelian strip with zero endpoint
vector, use instead the formula
\[
                  B_p(t)=(1-r)t+r\beta_{f_\xi}(t).
                                                               \tag{C.80.10}
\]
This last choice removes any dependence on
an interval ordering of the sphere parameters.
It is an increasing homeomorphism of $I$:
strict inequalities are preserved by its
convex combination, and both summands tend
to the same respective endpoints of $I$.
Its inverse is jointly continuous by strict
monotone bracketing. For each fixed $p,t$,
a parameter path from $*$ to $\xi$ shows by
the intermediate value theorem that an
original label in $f(\Sigma)\subset P$
realizes $B_p(t)$ at $t$. No continuous
choice of this path or label is used.

The other two prescriptions and their inverses
are jointly continuous by C.78's compact strict
gap and inverse-envelope arguments. All three
have exactly the relative conditions (C.80.7).
They and their inverses lie between the
corresponding compact $K$-evaluation envelopes.

Use the same countable comparison spaces as
in C.78--C.79: the joint character vector
space for an abelian strip; $I$ for a zero
endpoint vector; and each required scalar
plane with the actual action
\[
 \begin{split}
 T_e(k)(z,t)&=(z+c_e(k),\beta_k(t)),\\
 T_e(k)^{-1}(z,t)&=(z-c_e(k),\beta_k^{-1}(t)).
 \end{split}
                                                               \tag{C.80.11}
\]
For a rank-one vector choose a nonzero endpoint
character; the other endpoint character, when
present, is its actual real multiple. For a
rank-two vector use both planes, separately.
Their desired targets are exactly
$(r c_e(f_\xi),B_p(t))$ at source $(0,t)$.
Two extremal labels for the average have
the same full endpoint vector. A path in
$K$ joins them. The C.62 triangle in either
scalar plane, with real target $z$, realizes
the entire intervening fiber interval in
the single actual family $M_0=K K^{-1}K$.
The vector targets and (C.80.10) already
have labels in $P^5$ and $P$, respectively.
Thus $M_0$ is a common initial realizing
family for every listed comparison target.

All these actions of the actual Polish based
path group $\mathscr P(L)$ are transitive.
For the scalar planes this follows from the
transitive derived base image and the full
nonzero character image proved above; the
other cases use the actual vector image or
base orbit. The orbit map is endpoint action
on the named source point. Its stabilizer
is exactly the paths whose endpoint fixes
that point; the action kernel retains the
zero character and identity-base equations.
The verified Polish orbit theorem makes
these precise orbit maps open.

C.78's compact local pair-transport construction
therefore supplies, for every small identity
neighborhood, a compact family of actual maps
that realizes arbitrarily small pair transports
near each fixed comparison point. It is built
from compact preimages under these path-group
orbit maps and includes all countably many
spaces at once. Its proof uses no parameter
ordering. The second compact-product construction
gives one compact path-connected family $M\subset L$,
with nested $M_m$ starting at $M_0$, exactly
as in (C.78.13)--(C.78.14).

\emph{Relative fiber selection and the common-base splice.}
For each target equality just listed, take
the convex interval of $u_k(t)$ over labels
$k\in M_m$ satisfying that equality. These
intervals are nonempty. Small pair transports
correct an allowed label at $(p,t)$ to the
precise nearby target at $(p',t')$ and converge
to the identity. This is exactly the lower-limit
proof (C.78.16). Replace the intervals on
$D\times I$ by the values $0$, $u_{f_\xi}(t)$,
and $0$, in the order prescribed in (C.80.7).
The same relative approximation proof gives
continuous selections with these values exactly.

On a rank-two interval $I=(a,b)$ perform the
two scalar selections using the same $B_p$
and the same $M$. Denote them by $U_a,U_b$,
and use precisely the cutoff and blend of
(C.79.10):
\[
 \begin{split}
 \theta_I(t)&=\max\left(0,\min\left(1,
                     3\frac{t-a}{b-a}-1\right)\right),\\
 U&=(1-\theta_I)U_a+\theta_I U_b.
 \end{split}
                                                               \tag{C.80.12}
\]
All other intervals have their single selected
$U$. Define the maps and their actual inverses by
\[
 \begin{split}
 H_p(s,t)&=(s+U(p,t),B_p(t)),\\
 H_p^{-1}(s,v)&=
    (s-U(p,B_p^{-1}(v)),B_p^{-1}(v)).
 \end{split}
                                                               \tag{C.80.13}
\]
Every defining label in a fiber selection has
that identical actual base level. The rank-two
blend stays between the extrema of their union
at that level. The original source-plane C.62
triangle and its inverse bracketing therefore
give, uniformly over all parameters and strips,
\[
 H_p(x)\in M M^{-1}M(x),\qquad
 H_p^{-1}(x)\in M^{-1}M M^{-1}(x).
                                                               \tag{C.80.14}
\]
The required connecting paths lie in the
compact path-connected family $M$; no continuous
choice of them enters these set containments.

\emph{Every original component operator and every limit.}
Retain C.75's actual component generator
$q(s,t)=(\lambda s+\kappa,t+d)$ when it exists,
with $d>0$ and $\lambda\ne0$. Choose constructions
on interval representatives and transport them.
Endpoint characters multiply by $\lambda$;
the full-vector fibers are the same actual
label sets after that multiplication. Transport
their averages by those same labels, and their
fiber selections by $U'(p,t+d)=\lambda U(p,t)$.
All base prescriptions satisfy
$B'_p(t+d)=B_p(t)+d$. In particular (C.80.10)
has this identity by its displayed affine formula.
The cutoff in (C.80.12) has it as well.
For $I=\mathbb R$ stabilized by a component
translation, the actual line-bundle selection
of C.76--C.78 retains the equation with
$\lambda$, using its positive norm
$|\lambda|^{-t/d}|z|$. The relative approximation
works in that same bundle over $T$ times the
quotient base. Thus every operator equation,
including negative $\lambda$ and the unchanged
constant $\kappa$, holds exactly.

On each $a\in A$ put
$H_p(s,a)=(s+r u_{f_\xi}(a),a)$.
The finite-line limit proofs of C.78--C.79
are uniform for the new compact parameter
space $T$. In detail, compactness of $M$
and $f(\Sigma)$ gives uniform continuity
of their principal displacements on bounded
base intervals. Labels constrained at a
finite endpoint have precisely the specified
fraction there. For rank two, a fixed interval
approaching an endpoint uses the corresponding
selection because the cutoff is identically
zero or one nearby. For pairwise distinct
bounded intervals approaching a finite line,
both endpoints approach that line, by the
disjoint-interval argument of C.79. Each scalar
selection therefore has the same limiting
prescribed displacement; so does their blend.
Every arbitrary sequence has a subsequence
of one of those two forms. The compact inverse
base envelopes first give the same limiting
input line for (C.80.13), and the identical
argument then proves its inverse limit.
These assertions allow $r$ and $\xi$ to vary.

On a line over $A$, the continuous image
$\{u_{f_\eta}(a):\eta\in\Sigma\}$ is an interval
containing $0$ and $u_{f_\xi}(a)$. It therefore
contains $r u_{f_\xi}(a)$; the corresponding
inverse label realizes the negative displacement.
Thus the compact bounds (C.80.14) hold on
these lines too. Set $H_p=1$ on the original
boundary circle. Both word families there
are compact families of original closed-disk
maps fixing that circle. All their evaluations
at a point tending to a given circle point
converge uniformly to that same point. This
proves both required boundary limits, even
through varying strips and unbounded principal
coordinates.

The piecewise formulas preserve each base
interval and line and are inverse bijections.
The preceding limits make both maps jointly
continuous on $T\times\overline D$. Because
$\overline D$ is compact, this is continuity
$T\to\mathscr C_Q$ in its specified map-and-inverse
topology. The relative selections give exactly
the three equations in (C.80.2). Collapsing
$\{0\}\times S^n$ gives the cone on $S^n$,
homeomorphic to the closed $(n+1)$-disk, and
the basepoint ray remains the identity.
This is the asserted based filling.

The proof uses path connectedness and compactness
of $S^n$ for $n\geq1$. It makes no assertion
about $\pi_0$ or a single contraction of all
of $\mathscr C_Q$. It also supplies no homotopy
lift through the prime-end injection of C.73
for a different original compactification;
the remaining general C.12 extension question
retains that precise scope.
\end{proof}
